\section{Conclusions and Future Scope}
\label{sec:conclusion_future_scope}

This paper introduced \texttt{NyayaRAG}, a Retrieval-Augmented Generation framework tailored for realistic legal judgment prediction and explanation in the Indian common law system. By combining factual case details with retrieved statutory provisions and relevant precedents, our approach mirrors judicial reasoning more closely than prior methods that rely solely on the case text.
% 
Empirical results across prediction and explanation tasks confirm that structured legal retrieval enhances both outcome accuracy and interpretability. Pipelines enriched with statutes and precedents consistently outperformed baselines, as validated by lexical, semantic, and LLM-based (G-Eval) metrics, as well as expert feedback.

Future directions include extending to hierarchical verdict structures, integrating symbolic or graph-based retrieval, modeling temporal precedent evolution, and leveraging human-in-the-loop mechanisms. \texttt{NyayaRAG} marks a step toward court-aligned, explainable legal AI and sets the foundation for future research in retrieval-enhanced legal systems within under-represented jurisdictions.


% \section{Conclusion and Future Scope}
% \label{sec:conclusion_future_scope}
% In this paper, we presented \texttt{NyayaRAG}, a Retrieval-Augmented Generation (RAG) framework designed to simulate realistic legal judgment prediction and explanation within the Indian common law system. Unlike prior models that rely solely on the content of the current case, our approach integrates three key components, factual descriptions, relevant statutory provisions, and semantically similar prior cases, closely emulating the way judges reason in real courtroom settings.

% Through extensive experiments, we demonstrated that incorporating structured legal knowledge via RAG significantly improves both the predictive performance and the legal soundness of generated explanations. Specifically, pipelines augmented with statutes and precedents yielded higher accuracy and interpretability, as confirmed by both traditional NLP metrics and LLM-based evaluation using G-Eval. Our work underscores the importance of simulating actual legal reasoning processes to develop more faithful and trustworthy AI systems for legal applications.

% Future work may explore several promising directions. First, we aim to extend our framework to multi-class or hierarchical verdict structures that better capture real-world legal complexity. Second, while our current retrieval is based on dense semantic similarity, future iterations could integrate symbolic reasoning or graph-based legal knowledge for more structured retrieval. Third, we plan to incorporate temporal dynamics of precedent evolution, enabling the system to weigh older versus newer case law appropriately. Lastly, we envision incorporating human-in-the-loop feedback and expert validation to further align AI predictions with judicial expectations.

% % \noindent
% By aligning AI with the foundational principles of Indian jurisprudence, \texttt{NyayaRAG} contributes toward the long-term vision of explainable, realistic, and accessible legal AI. We hope this work sparks further research into retrieval-enhanced, court-aligned AI systems in underrepresented legal ecosystems.

